$$
	y = \beta_0 + \beta_1 x_1 + \beta_2 x_2 + ... + \beta_n x_n + \epsilon
$$

\par
For each \textit{feature} $x$ there is an adjustable parameter, $\beta_1$,
which corresponds to its slope or gradient (often called the \textbf{weight})
as well as a intersection term $\beta_0$ (often called the \textbf{bias}).

\par
This is a \textit{parametric} model; here we assume the relation between the
dependent and the independent variables is linear, and the noise is from a
(\textbf{Gaussian}) normal distribution.

\par
Our data especially when working with tabular data, is usually stored in a
\texttt{DataFrame} (\texttt{df}). One can you use the \texttt{pandas} library to
handle this data structure, or convert it to a \texttt{numpy} array.

\par
There exists a closed form solution for this problem:

$$
	y = \beta_1x+\beta_0
$$

The OLS equation for $\beta_1$:

$$
	\beta_1 = \frac{
		\sum((x_i - \overline{x})(y_i - \overline{y}))}{
		\sum(x_i - \overline{x})^2
	}
$$

And for $\beta_0$:

$$
	\beta_0 = \overline{y} - \beta_1 \overline{x}
$$

\par
Our \texttt{DataFrame} is basically an augmented coefficient matrix:

% TODO: Include example

\par
We can treat our \texttt{DataFrame} as if it were a system of linear equations
where values of $x_{mn}$ represent the coefficients:

%TODO: Include example matrix
$$
	%	\left\{ \begin{array}{l}
	%	y_1 = \beta_0 + \beta_1 x_{11} + \beta_2 x_{12} + ... + \beta_n x_{1n} \\
	%	y_2 = \beta_0 + \beta_1 x_{21} + \beta_2 x_{22} + ... + \beta_n x_{2n} \\
	%	... \\
	%	y_m = \beta_0 + \beta_1 x_{m1} + \beta_2 x_{m2} + ... + \beta_n x_{mn} \\
	%	\right\}
$$

\subsection{Loss function}

\par
There are 2 options in calculating the loss function for linear regression:

\begin{itemize}
	\item \textbf{error:} Distance between points (not used) $y_i-\hat{y_i}$
	\item \textbf{Absolute distance:} $L1 loss = |y_i-\hat{y_i}|$
	\item \textbf{The square:} $L2 loss = (y_i-\hat{y_i})^2$
\end{itemize}

\par
A correct loss function should satisfy these requirements:

% TODO:

Each function has a different behavior:

\begin{itemize}
	\item $L1 loss$ is not very sensitive to \textit{outliers}
	\item $L2 loss$ is very sensitive to \textit{outliers}
\end{itemize}

In code:

\begin{lstlisting}[language=Python, caption=Loss functions]
# y
y_true

# y hat
y_pred

# L1 loss = AE (absolute error)
abs(y_true - y_pred)

# L1 cost = MAE (mean absolute error)
sum(abs(y_true - y_pred)) / len(y_true)

# L2 loss = SE (squared error)
(y_true - y_pred) ** 2

# L2 cost = MSE (mean squared error)
sum((y_true - y_pred) ** 2) / len(y_true)
\end{lstlisting}

\par
We want to minimize the cost or \textit{objective} function $J$. For example, in
linear regression, we use ordinary least squares, in other words the $L2 loss$
and the $L2 cost$.

% TODO:
%\begin{array}
%	J(\beta) & = MSE
%	         & = \frac{1}{m} \sum_{i=1}^{m} (y_i - \hat{y_i})^2
%\end{array}

\par
The \texttt{argmin(J)} function gives us the optial values of  $\beta_1$ and
$\beta_0$.

\par
$J$ is piecewise-continuous (for $L1 loss$)and the minima is multivalued for this
example, which is ambiguous (not great). For $J$ in $L2 loss$ is smooth, it is a
convex function which means it has a unique global minima. In a 2 dimensional
problem, $J$ takes a parabolic surface as a form, having the same dimension as
the number of parameters.

\subsection{Gradient Descent}

\par
The gradient ($G$) of a surface is given by:

$$
	G = \frac{\partial J(\beta)}{\partial \beta_0}
$$

where $\beta$ is a vector of parameters.

% TODO: more theory on this

\par
Gradient descent does not know the form of the parabola, meaning that if $\eta$
is not set appropriately, the step can be too big and can quickly diverge away.

% TODO: Include graphs of $\eta$too small, good or too big

\par
The Aagrad Gradient Algorithm changes the $\eta$ on the fly

% TODO: Equation for $\eta$

\par
What happens if the features have very different scales? The steps taken in both
$\beta$ values are very different. This can be solved by finding the maximum and
minimum values of each feature and normalizing them to a common scale. This
makes it so that the previous cost function $J$ for the $L2 loss$ is now a more
circular paraboloid. This is called \textbf{feature scaling} or
\textbf{normalization}.

\par
With parametric estimators the solvers have better numerical stability if the
data is \textit{normalized}. Here are some linear transformations.

\begin{itemize}
	\item \texttt{sklearn.preprocessing.MinMaxScaler()}: $[0,1]$
	\item \texttt{sklearn.preprocessing.StandardScaler()}: $\eta = 0, \sigma = 1$
	\item \texttt{sklearn.preprocessing.Normalizer()}: in case of outliers
\end{itemize}

\par
Usage example:

\begin{lstlisting}[language=Python, caption=Feature scaling]
from sklearn.preprocessing import RobustScaler

RS = RobustScaler()

X_train = RS.fit_transform(X_train)
X_test = RS.transform(X_test)
\end{lstlisting}

An example of a real neural network loss surface:

% TODO: Include example

\par
The reason why numerical methods are used is because there aren't any methods
that analytically find a global minimum in these real neural network loss
surface models. This makes them very costly to train.



